\chapter*{CALCULOS JUSTIFICATIVOS}
\addcontentsline{toc}{chapter}{CALCULOS JUSTIFICATIVOS}

% Reinicio de contadores para numeracion coherente con el capitulo 2.
\setcounter{section}{0}
\renewcommand{\thesection}{2.\arabic{section}.}
\renewcommand{\thesubsection}{2.\arabic{section}.\arabic{subsection}.}
\renewcommand{\thesubsubsection}{2.\arabic{section}.\arabic{subsection}.\arabic{subsubsection}.}
\renewcommand{\theHsection}{calculos.\arabic{section}}
\renewcommand{\theHsubsection}{calculos.\arabic{section}.\arabic{subsection}}
\renewcommand{\theHsubsubsection}{calculos.\arabic{section}.\arabic{subsection}.\arabic{subsubsection}}

\section{NORMAS APLICABLES Y BASES DE DISENO}

Los calculos justificativos se desarrollan para una instalacion residencial monofasica en baja tension. Las normas consideradas como base son el Codigo Nacional de Electricidad - Utilizacion, la Norma Tecnica EM.010 del Reglamento Nacional de Edificaciones, la Norma DGE de Simbolos Graficos en Electricidad y los criterios tecnicos de la empresa concesionaria local.

La auditoria normativa del capitulo 2 corrigio tres criterios importantes:
\begin{itemize}
  \item La Regla CNE-U 050-204 no se usa como sustento para el factor \(1.25\), porque la revision del texto disponible la relaciona con escuelas y no con vivienda.
  \item El factor \(1.25\) se mantiene solo como criterio preliminar conservador para verificar la capacidad del conductor. Su sustento normativo final debe revisarse en el CNE-U antes del cierre.
  \item Los circuitos derivados de alumbrado se dimensionan con conductor minimo de \(2.5\text{ mm}^2\), conforme a la observacion de auditoria sobre CNE-U 030-002. No se adopta \(1.5\text{ mm}^2\) como seccion definitiva para circuitos derivados de alumbrado.
\end{itemize}

\section{PARAMETROS ELECTRICOS DEL SISTEMA}

Para el calculo preliminar se adoptan los parametros de la Tabla \ref{tab:parametros_calculo}.

\begin{table}[H]
\centering
\caption{Parametros electricos usados en el calculo preliminar.}
\label{tab:parametros_calculo}
\begin{tabularx}{\textwidth}{l c L}
\toprule
\textbf{Parametro} & \textbf{Valor} & \textbf{Observacion} \\
\midrule
Tension nominal & \(220\text{ V}\) & Suministro monofasico residencial por confirmar con concesionaria. \\
Frecuencia & \(60\text{ Hz}\) & Valor usual del sistema electrico nacional. \\
Factor de potencia & \(0.90\) & Valor preliminar para cargas residenciales combinadas. \\
Resistividad del cobre & \(0.0178\ \Omega\cdot\text{mm}^2/\text{m}\) & Valor unico usado en motor, tablas y formulas. \\
Factor de diseno del conductor & \(1.25\) & Criterio conservador preliminar; REVISAR sustento normativo final. \\
Limite de caida alimentador & \(1.5\%\) & Criterio de diseno conservador; el limite normativo debe verificarse en CNE-U. \\
Limite de caida derivados & \(2.5\%\) & Criterio usado para circuitos derivados. \\
\bottomrule
\end{tabularx}
\end{table}

\section{DATOS BASE DE AREAS}

Las areas arquitectonicas disponibles aun no son definitivas. La memoria descriptiva anterior, el cuestionario y los croquis no coinciden plenamente, por lo que se trabaja con un escenario de calculo provisional y se marca el estado de cada dato.

\begin{table}[H]
\centering
\caption{Datos base de areas usados para calculo preliminar.}
\label{tab:areas_base_calculo}
\begin{small}
\begin{tabularx}{\textwidth}{l c l L}
\toprule
\textbf{Dato} & \textbf{Valor (m$^2$)} & \textbf{Estado} & \textbf{Fuente / observacion} \\
\midrule
Area de terreno & 134.18 & preliminar & respuestas-cuestionario-aquiles.md \\
Area techada primer piso & Por confirmar & por confirmar & croquis y planos CAD aun no acotados \\
Area techada segundo piso & 42.56 & preliminar & respuestas-cuestionario-aquiles.md \\
Area techada total de calculo & 120.00 & preliminar & escenario provisional de calculo; debe validarse con Aquiles \\
\bottomrule
\end{tabularx}
\end{small}
\end{table}

El area techada total usada para el escenario preliminar es \(120.00\text{ m}^2\). Este valor no reemplaza la medicion final de planos; solo permite evaluar la demanda basica por area mientras Aquiles confirma las medidas reales por piso.

\section{LEVANTAMIENTO PRELIMINAR DE CARGAS}

El levantamiento de cargas se organiza por circuitos derivados. Las potencias unitarias y cantidades provienen de las respuestas preliminares, los croquis de distribucion y el cuadro de cargas generado por el motor Python. El circuito C6 se trata como dato sensible porque la cocina puede ser a gas o electrica.

Para evitar cerrar una condicion no confirmada, el motor calcula dos escenarios:
\begin{itemize}
  \item \textbf{Escenario 1:} cocina a gas o sin cocina electrica de alta potencia.
  \item \textbf{Escenario 2:} cocina electrica considerada como caso conservador preliminar.
\end{itemize}

\begin{table}[H]
\centering
\caption{Comparacion de escenarios de maxima demanda preliminar.}
\label{tab:escenarios_demanda}
\begin{scriptsize}
\resizebox{\textwidth}{!}{%
\begin{tabular}{l c c c c c c c c}
\toprule
\textbf{Escenario} & \textbf{P.I. (W)} & \textbf{MD circ. (W)} & \textbf{MD CNE (W)} & \textbf{MD adopt. (W)} & \textbf{Ib (A)} & \textbf{Cond. (mm$^2$)} & \textbf{ITM} & \textbf{dV alim. (\%)} \\
\midrule
Escenario 1 - cocina a gas o sin cocina electrica de alta potencia & 9,058 & 7,078 & 3,500 & 7,078 & 35.75 & 10.0 & 2P-40A & 0.694 \\
Escenario 2 - cocina electrica considerada & 11,658 & 10,358 & 9,500 & 10,358 & 52.31 & 16.0 & 2P-63A & 0.635 \\
\bottomrule
\end{tabular}%
}
\end{scriptsize}
\end{table}

\section{POTENCIA INSTALADA}

La potencia instalada se calcula con la suma de las potencias nominales de las salidas y equipos considerados:
\begin{equation}
P_{inst}=\sum_{j=1}^{n} \left(N_j \times P_j\right)
\end{equation}

Donde \(N_j\) es la cantidad de elementos de la carga tipo \(j\) y \(P_j\) es la potencia nominal unitaria en watts. En el escenario sin cocina electrica de alta potencia se obtiene \(P_{inst}=9,058\text{ W}\). En el escenario con cocina electrica considerada se obtiene \(P_{inst}=11,658\text{ W}\).

\section{MAXIMA DEMANDA}

La maxima demanda se verifica por dos criterios y se adopta el mayor resultado para no subdimensionar el alimentador:
\begin{equation}
P_{MD,adoptada}=\max\left(P_{MD,circuitos},P_{MD,CNE}\right)
\end{equation}

El metodo por circuitos calcula:
\begin{equation}
P_{MD,circuitos}=\sum \left(P_{inst,circuito}\times FD_{circuito}\right)
\end{equation}

El metodo de carga basica de vivienda, segun el criterio CNE-U 050-200 revisado en la auditoria, considera:
\begin{equation}
P_{base}=2,500\text{ W}+1,000\text{ W}\times
\left\lceil\frac{\max(A_{techada}-90,0)}{90}\right\rceil
\end{equation}

Para \(A_{techada}=120\text{ m}^2\), la carga basica es:
\begin{equation}
P_{base}=2,500\text{ W}+1,000\text{ W}=3,500\text{ W}
\end{equation}

Si la vivienda tuviera cocina electrica, el escenario conservador incorpora \(6,000\text{ W}\) para esa carga, obteniendo \(P_{MD,CNE}=9,500\text{ W}\). El metodo por circuitos del mismo escenario arroja \(10,358\text{ W}\); por lo tanto, el valor adoptado provisionalmente para dimensionamiento conservador es:
\begin{equation}
P_{MD,adoptada}=10,358\text{ W}
\end{equation}

Este valor no confirma que la vivienda tendra cocina electrica. Se usa para evitar que el alimentador y las protecciones queden subdimensionados mientras Aquiles confirma el tipo real de cocina.

\section{CORRIENTE DE EMPLEO Y CORRIENTE DE DISENO}

La corriente de empleo del alimentador se calcula como:
\begin{equation}
I_b=\frac{P_{MD,adoptada}}{V\times\cos\phi}
\end{equation}

Para el escenario conservador preliminar:
\begin{equation}
I_b=\frac{10,358\text{ W}}{220\text{ V}\times0.90}=52.31\text{ A}
\end{equation}

La corriente de diseno del conductor se verifica con un margen conservador:
\begin{equation}
I_d=1.25\times I_b=1.25\times52.31\text{ A}=65.39\text{ A}
\end{equation}

El factor \(1.25\) no queda citado como exigencia de la Regla CNE-U 050-204. En esta version se conserva como criterio preliminar de diseno y se mantiene marcado para revision normativa final.

\section{SELECCION DEL ALIMENTADOR PRINCIPAL}

La seleccion del alimentador se realiza con dos controles:
\begin{itemize}
  \item Capacidad del conductor: \(I_d \leq I_z\).
  \item Coordinacion de proteccion: \(I_b \leq I_{ITM} \leq I_z\).
\end{itemize}

Con el escenario conservador preliminar, el motor selecciona conductor de cobre de \(16.0\text{ mm}^2\), con ampacidad referencial \(I_z=66\text{ A}\), e interruptor termomagnetico bipolar de \(63\text{ A}\). La verificacion principal queda:
\begin{equation}
65.39\text{ A}\leq66\text{ A}
\end{equation}
\begin{equation}
52.31\text{ A}\leq63\text{ A}\leq66\text{ A}
\end{equation}

Si Aquiles confirma cocina a gas y no existe otra carga electrica de alta potencia, el escenario alternativo permite revisar el alimentador a \(10.0\text{ mm}^2\) e interruptor general bipolar de \(40\text{ A}\). Esa opcion no debe cerrarse hasta confirmar las cargas reales.

\section{CIRCUITOS DERIVADOS}

La distribucion preliminar considera ocho circuitos derivados:
\begin{itemize}
  \item C1: Alumbrado del primer piso.
  \item C2: Tomacorrientes generales del primer piso.
  \item C3: Tomacorriente auxiliar de cocina del primer piso.
  \item C4: Alumbrado del segundo piso.
  \item C5: Tomacorrientes generales del segundo piso.
  \item C6: Cocina del segundo piso o carga de cocina de alta potencia, por confirmar.
  \item C7: Lavadora del segundo piso.
  \item C8: Bomba de agua exterior.
\end{itemize}

Los circuitos C1 y C4 usan \(2.5\text{ mm}^2\) como seccion minima para alumbrado derivado. El circuito C6 se dimensiona de manera conservadora con \(10.0\text{ mm}^2\) en el escenario con cocina electrica preliminar.

\section{VERIFICACION DE CAIDA DE TENSION}

La caida de tension monofasica se calcula con:
\begin{equation}
\Delta V=\frac{2\times L\times I_b\times\rho}{S}
\end{equation}

Y el porcentaje se obtiene mediante:
\begin{equation}
\%\Delta V=\frac{\Delta V}{V}\times100
\end{equation}

Para el alimentador del escenario conservador preliminar, con recorrido supuesto de \(12\text{ m}\), conductor de \(16.0\text{ mm}^2\), \(\rho=0.0178\ \Omega\cdot\text{mm}^2/\text{m}\) e \(I_b=52.31\text{ A}\), la caida de tension calculada es \(0.635\%\). Este resultado cumple el criterio conservador adoptado de \(1.5\%\), aunque la longitud real medidor--TG debe confirmarse.

\section{CUADRO RESUMEN DE CARGAS}

El cuadro siguiente corresponde al escenario de dimensionamiento conservador preliminar con cocina electrica considerada. No confirma que la cocina electrica vaya a instalarse; solo documenta la hipotesis que gobierna el dimensionamiento actual.

\begin{table}[H]
\centering
\caption{Cuadro de cargas preliminar - Escenario 2 - cocina electrica considerada.}
\label{tab:cuadro-cargas-cocina-electrica-preliminar}
\begin{small}
\begin{tabular}{l l c c c}
\toprule
\textbf{Circuito} & \textbf{Descripcion (Uso)} & \textbf{P.I. (W)} & \textbf{F.D.} & \textbf{M.D. (W)} \\
\midrule
C1 & Alumbrado 1er Piso & 80 & 1.00 & 80.0 \\
C2 & Tomacorrientes 1er Piso & 1,260 & 0.70 & 882.0 \\
C3 & Cocina 1er Piso (Auxiliar) & 300 & 0.80 & 240.0 \\
C4 & Alumbrado 2do Piso & 132 & 1.00 & 132.0 \\
C5 & Tomacorrientes 2do Piso & 2,340 & 0.70 & 1,638.0 \\
C6 & Cocina electrica 2do Piso (escenario por confirmar) & 6,000 & 1.00 & 6,000.0 \\
C7 & Lavadora 2do Piso & 800 & 0.80 & 640.0 \\
C8 & Bomba de Agua Exterior & 746 & 1.00 & 746.0 \\
\midrule
\multicolumn{2}{l}{\textbf{Total instalado / MD circuitos}} & \textbf{11,658} & & \textbf{10,358.0} \\
\bottomrule
\end{tabular}
\end{small}
\end{table}

\section{CUADRO RESUMEN DE CONDUCTORES, TUBERIAS Y PROTECCIONES}

La Tabla \ref{tab:conductores-protecciones-cocina-electrica-preliminar} muestra la verificacion de conductor, proteccion y caida de tension generada por el motor Python. En la columna de estado, las tres verificaciones corresponden a conductor, coordinacion e indice de caida de tension.

\begin{table}[H]
\centering
\caption{Seleccion preliminar de conductores y protecciones - Escenario 2 - cocina electrica considerada.}
\label{tab:conductores-protecciones-cocina-electrica-preliminar}
\begin{scriptsize}
\resizebox{\textwidth}{!}{%
\begin{tabular}{l c c c c c l c c c}
\toprule
\textbf{Circ.} & \textbf{P.I. (W)} & \textbf{M.D. (W)} & \textbf{Ib (A)} & \textbf{Id (A)} & \textbf{Cond.} & \textbf{ITM} & \textbf{Iz (A)} & \textbf{dV (\%)} & \textbf{Estado} \\
\midrule
Alim. & 11,658 & 10,358 & 52.31 & 65.39 & 16.0 mm$^2$ & 2P-63A & 66 & 0.635 & CUMPLE \\
\midrule
C1 & 80 & 80.0 & 0.40 & 0.51 & 2.5 mm$^2$ & 2P-10A & 20 & 0.039 & CUMPLE/CUMPLE/CUMPLE \\
C2 & 1,260 & 882.0 & 4.46 & 5.57 & 2.5 mm$^2$ & 2P-10A & 20 & 0.519 & CUMPLE/CUMPLE/CUMPLE \\
C3 & 300 & 240.0 & 1.21 & 1.51 & 2.5 mm$^2$ & 2P-10A & 20 & 0.078 & CUMPLE/CUMPLE/CUMPLE \\
C4 & 132 & 132.0 & 0.67 & 0.83 & 2.5 mm$^2$ & 2P-10A & 20 & 0.095 & CUMPLE/CUMPLE/CUMPLE \\
C5 & 2,340 & 1,638.0 & 8.27 & 10.34 & 2.5 mm$^2$ & 2P-10A & 20 & 1.339 & CUMPLE/CUMPLE/CUMPLE \\
C6 & 6,000 & 6,000.0 & 30.30 & 37.88 & 10.0 mm$^2$ & 2P-32A & 50 & 0.981 & CUMPLE/CUMPLE/CUMPLE \\
C7 & 800 & 640.0 & 3.23 & 4.04 & 2.5 mm$^2$ & 2P-10A & 20 & 0.314 & CUMPLE/CUMPLE/CUMPLE \\
C8 & 746 & 746.0 & 3.77 & 4.71 & 2.5 mm$^2$ & 2P-10A & 20 & 0.341 & CUMPLE/CUMPLE/CUMPLE \\
\bottomrule
\end{tabular}%
}
\end{scriptsize}
\end{table}

\section{SISTEMA DE PUESTA A TIERRA}

El proyecto contempla un sistema de puesta a tierra residencial asociado al tablero general y a los tomacorrientes con contacto de tierra. La resistencia objetivo preliminar se mantiene como dato por confirmar, porque debe validarse con la normativa aplicable, la disponibilidad de terreno natural y la medicion de resistividad o resistencia en obra.

El calculo detallado del pozo de tierra no se cierra en esta version. Deben confirmarse ubicacion, tipo de electrodo, condiciones del suelo y valor objetivo antes de aprobar el dimensionamiento final del S.P.A.T.

\section{OBSERVACIONES TECNICAS}

\begin{itemize}
  \item Los resultados del presente capitulo son preliminares y estan sujetos a validacion de areas, cargas reales y recorridos de canalizacion.
  \item El escenario recomendado para no subdimensionar, mientras no haya respuesta de Aquiles, es el escenario conservador con cocina electrica considerada.
  \item Si Aquiles confirma cocina a gas, ausencia de ducha electrica y ausencia de terma electrica, se podra recalcular y posiblemente reducir alimentador e interruptor general.
  \item La seleccion de \(2.5\text{ mm}^2\) para alumbrado corrige la observacion critica de la auditoria.
  \item La cita a CNE-U 050-204 fue retirada como fundamento del factor \(1.25\).
  \item Las longitudes usadas en caida de tension son supuestas; deben corregirse con el plano electrico y la ubicacion real del medidor y tablero general.
\end{itemize}
